\section{Introduction}
Federated learning (FL) has become a standard paradigm for collaborative model training without centralizing raw data, with FedAvg~\cite{mcmahan2017communication} and heterogeneity-aware extensions such as FedProx~\cite{li2020federated} establishing the modern optimization framework for decentralized learning under communication and systems constraints. Yet, even in this mature line of work, convergence remains fragile when clients are statistically heterogeneous, intermittently available, or resource-constrained, because the global objective is optimized through delayed and imperfect local surrogates rather than through direct centralized gradients. This tension is especially acute when the model is too large or too expensive to train fully on-device, motivating alternatives that reduce client-side memory, computation, and communication while preserving privacy. 

Split learning (SL) offers one such alternative by partitioning a neural network across clients and servers so that clients compute only shallow layers while deeper layers are executed remotely~\cite{vepakomma2018split}. Hierarchical client-edge-cloud variants extend this idea to multi-tier deployments, improving scalability and reducing end-to-end latency through intermediate aggregation~\cite{liu2020client}. However, the resulting optimization remains tightly coupled across tiers: standard split methods depend on cross-tier forward/backward message exchange, while hierarchical FL variants largely analyze periodic model aggregation rather than decoupled representation-level coordination. As a result, these methods reduce device burden, but they do not directly explain how learning behaves when upper-tier training is driven by partial, delayed, or asynchronous summaries of lower-tier representations. 

A second line of work replaces gradient exchange with more abstract knowledge transfer. Prototype-based FL communicates class-level representations rather than raw parameters, as in FedProto~\cite{tan2022fedproto}, while heterogeneous-model fusion methods such as FedDF~\cite{lin2020ensemble} and generator-based approaches such as FedGen~\cite{zhu2021data} train on unlabeled or synthesized knowledge at the server. These studies are highly relevant because they move coordination from parameter space into feature or distribution space, which is precisely the direction needed for communication-efficient upper-tier training. Nevertheless, existing prototype and generative FL methods are mostly developed for flat server-client topologies or for server-side regularization/distillation; they do not study hierarchical split training in which edge and cloud models are optimized over prototype-induced synthetic distributions that may themselves be incomplete, stale, and temporally inconsistent. 

This limitation becomes more severe under asynchronous and serverless-style execution. Prior asynchronous FL analyses, including Asynchronous Federated Optimization~\cite{xie2019asynchronous}, FedBuff~\cite{nguyen2022federated}, and later general-theory treatments of heterogeneous delayed updates~\cite{fraboni2023general}, study stale model updates and delayed aggregation in parameter space. In parallel, serverless FL platforms such as FedLess~\cite{grafberger2021fedless} and LIFL~\cite{qi2024lifl} show that event-driven, elastic, and scale-to-zero execution can improve scalability and resource efficiency for FL systems. However, from the standpoint of hierarchical split-federated prototyping, the most important consequence of such execution is not resource allocation but learning instability: upper-tier workers are invoked episodically, see only partial arrivals of prototype packets, and may optimize on distributions synthesized from outdated lower-tier statistics. In this setting, the central object of analysis is no longer a buffered stream of model updates but a sequence of shifting prototype-induced measures. 

This paper addresses that gap through Episodic Hierarchical Split-Federated Prototyping (E-HSFP), a convergence-guided learning framework. Our starting point is to treat serverless-style episodes as a distributional learning problem: edge and cloud models are trained on synthetic feature distributions induced by partial and delayed prototype summaries, and convergence depends on whether the resulting drift can be bounded. E-HSFP therefore introduces five learning mechanisms, each tied to a specific source of drift or estimation error. Episodic prototype memory stabilizes temporal variation by replaying historical class-wise statistics, reducing abrupt shifts in the effective training distribution. Reliability-aware prototype aggregation down-weights stale, noisy, or weakly supported prototypes, thereby reducing aggregation bias under partial arrivals. Prototype dropout regularization trains upper-tier models against missing prototype events and incomplete episodes. Prototype replay consistency aligns current and memory-induced synthetic features so that upper-tier representations do not diverge across episodes. Finally, a residual prototype generator enriches diagonal-Gaussian feature generation and reduces synthetic-distribution approximation error.

The main theoretical goal is to prove that E-HSFP converges to a stationary neighborhood under non-convex optimization with explicit dependence on both classical and episodic error sources. Concretely, the target bound decomposes the stationarity gap into stochastic gradient noise, periodic sub-model aggregation divergence, partial prototype updates, stale prototype updates, memory replay error, and synthetic generation error. This decomposition is important not only for analysis but also for design: each E-HSFP component is introduced because it shrinks one or more identifiable terms in the bound, turning serverless-style robustness motivation into a convergence-guided algorithmic framework.

The empirical program is designed to mirror that theory. Beyond final accuracy, the planned evaluation measures convergence curves, gradient-norm trajectories, and prototype-induced distributional drift through distances such as Wasserstein or MMD across rounds. It stress-tests the method under missing prototype packets and stale prototype updates, compares it against standard FL, split/hierarchical FL, and prototype/generative baselines, and performs component ablations for episodic prototype memory, reliability-aware prototype aggregation, prototype replay consistency, and the residual prototype generator. To test generality beyond classification, the protocol also includes a dense-prediction setting, e.g., medical image segmentation, where upper-tier drift can be especially damaging to stable optimization.

In summary, this work makes four contributions:
\begin{itemize}
    \item We formulate episodic hierarchical split-federated prototyping as optimization over time-varying prototype-induced feature distributions.
    \item We propose E-HSFP, a learning-based framework that counteracts partial prototype updates, stale prototype updates, and synthetic-distribution drift through episodic prototype memory, reliability-aware prototype aggregation, prototype dropout regularization, prototype replay consistency, and a residual prototype generator.
    \item We establish a convergence guarantee to a stationary neighborhood with an explicit decomposition of stochastic, aggregation, and episodic drift errors.
    \item We specify a convergence-centric empirical protocol for measuring stability, robustness, and generality rather than relying on endpoint accuracy alone. Under this protocol, E-HSFP consistently improves on H-SFP and is robust to prototype dropout and serverless-style disruptions, although it does not close the gap to the strongest weight-sharing baselines.
\end{itemize} 
